\section{Metric-response theorem from the reduced BHSM operator}
\label{sec:bhsm-metric-response}

The preceding section identifies the physical $n=2$ topographic carrier and
its reduced phase-space pair $X_2=(q_2,\Pi_2)^T$.  The next required bridge is
from that mode to the metric perturbation seen by late-time observables.  This
bridge follows directly from the mixed physical Hessian and does not require
a second independently propagating scalar.

\subsection{Coupled quadratic block}

On the gauge-fixed physical tangent space, let $h$ denote the retained scalar
metric perturbation variables and let $X_2$ denote the physical topographic
mode data.  After first-order reduction of any derivative mixing, the linear
metric equation can be written
\begin{equation}
\mathcal A_g h
+
\mathcal B_{g2}X_2
+
J_g
=
0,
\label{eq:metric-block}
\end{equation}
where $\mathcal A_g$ is the physical metric Hessian block,
$\mathcal B_{g2}$ is the metric--topographic mixed block, and $J_g$ is any
independent matter/environment source.

The BHSM eta--metric mixed second variation supplies a concrete microscopic
source for such a mixed block.  In the notation used earlier in the
manuscript, that second variation contains terms of the schematic form
\begin{equation}
B_{g\phi}[\gamma,\phi]
=
\int d\mu_g\,w\,\gamma^{\mu\nu}
\left[
2F''(X)(j\!\cdot\! d\phi)A_{\mu\nu}
+
2F'(X)j_{(\mu}\partial_{\nu)}\phi
-
F'(X)g_{\mu\nu}(j\!\cdot\!d\phi)
\right].
\label{eq:bhsm-mixed-source}
\end{equation}
Equation~\eqref{eq:bhsm-mixed-source} is not asserted to be the complete
cosmological propagator.  It establishes that the BHSM action already owns a
metric--topographic mixing channel from which the reduced response map can be
constructed.

\subsection{Induced-metric theorem}

\paragraph{Theorem 2 (induced metric response).}
Assume that $\mathcal A_g$ is invertible on the chosen physical metric domain.
Then the metric response induced by the physical $n=2$ mode is unique and is
given by
\begin{equation}
\boxed{
h_2
=
h_{\rm src}
+
\mathcal R_{g2}X_2,
\qquad
\mathcal R_{g2}
=
-\mathcal A_g^{-1}\mathcal B_{g2},
}
\label{eq:metric-response-map}
\end{equation}
with
\begin{equation}
h_{\rm src}=-\mathcal A_g^{-1}J_g.
\end{equation}
No additional physical scalar degree of freedom is introduced.

\paragraph{Proof.}
Equation~\eqref{eq:metric-block} is linear on the physical tangent space.
Invertibility of $\mathcal A_g$ therefore gives
\[
h
=
-\mathcal A_g^{-1}J_g
-
\mathcal A_g^{-1}\mathcal B_{g2}X_2.
\]
The second term is the unique response sourced by the retained topographic
mode.  Since $X_2$ already contains the complete reduced canonical data of
that mode, the induced metric response carries no additional independent
initial datum. \hfill$\square$

\subsection{Schur-complement consistency}

In second-order notation, before first-order phase-space reduction, the
quadratic action can be represented by a block operator
\begin{equation}
\mathbb H_2
=
\begin{pmatrix}
\mathcal A_g & \mathcal B_2\\
\mathcal B_2^\dagger & \mathcal C_2
\end{pmatrix}.
\end{equation}
Eliminating the metric block gives
\begin{equation}
\boxed{
\mathcal K_2
=
\mathcal C_2
-
\mathcal B_2^\dagger
\mathcal A_g^{-1}
\mathcal B_2 ,
}
\label{eq:metric-schur}
\end{equation}
which is the same reduced operator that governs the action-owned mode.
Thus the metric-response map and the reduced topographic evolution are two
faces of the same Schur/Feshbach reduction, rather than independent model
assumptions.

\subsection{Closed-$S^3$ harmonic preservation}

For an isotropic closed background, $\mathcal A_g$ and the action-owned mixed
block commute with the background isometry action.  Consequently the
response operator $\mathcal R_{g2}$ preserves the $n=2$ representation unless
an additional anisotropic source is inserted.

\paragraph{Corollary 2.1.}
For the pure $A_{4i}$ component of the physical $n=2$ mode, every scalar
component of the induced metric response has the observer-sky factor
\begin{equation}
\hat{\mathbf n}\!\cdot\!\hat{\mathbf p}.
\end{equation}
Hence one may write
\begin{align}
\Phi_2(z,\hat{\mathbf n})
&=
\Phi_{2,\rm amp}(z)
(\hat{\mathbf n}\!\cdot\!\hat{\mathbf p}),
\\
\Psi_2(z,\hat{\mathbf n})
&=
\Psi_{2,\rm amp}(z)
(\hat{\mathbf n}\!\cdot\!\hat{\mathbf p}),
\label{eq:metric-potential-factorization}
\end{align}
with amplitudes determined by
Eq.~\eqref{eq:metric-response-map}.

\subsection{Observable response operator}

Let $\mathcal L_{\mathcal O,z}$ denote the standard linear observation
functional for a scalar observable $\mathcal O$ at redshift $z$.  Combining
the induced metric response with any matter response generated by the same
reduced mode gives
\begin{equation}
\delta\mathcal O_2
=
\mathcal L_{\mathcal O,z}
\left[
\mathcal R_{g2}X_2
+
\mathcal R_{m2}X_2
\right].
\label{eq:observable-response-operator}
\end{equation}
The observable transfer amplitude is therefore
\begin{equation}
\boxed{
\mathcal A_{\mathcal O}(z)
=
\mathcal L_{\mathcal O,z}
\left[
\mathcal R_{g2}(z)
+
\mathcal R_{m2}(z)
\right]X_2(z),
}
\label{eq:observable-amplitude-from-bhsm}
\end{equation}
up to independent environmental sources explicitly separated from the
single-mode prediction.

Equation~\eqref{eq:observable-amplitude-from-bhsm} is the direct BHSM-native
bridge from the reduced $n=2$ phase-space state to the measurable amplitudes
used in the distance, expansion, growth and lensing sectors.

\subsection{Scientific closure condition}

The remaining numerical closure object is now sharply defined:
\begin{equation}
\boxed{
\mathcal R_{g2}(z)
=
-\mathcal A_g^{-1}(z)\mathcal B_{g2}(z)
}
\end{equation}
together with the matter-response block $\mathcal R_{m2}(z)$ on the frozen
reference branch.  Once these operators are evaluated from the same reduced
action, the redshift-dependent observable kernels cease to be independent
phenomenological inputs.

A curved four-dimensional EFT realization is accepted as an equivalent
representation only if its physical response map agrees with
Eq.~\eqref{eq:metric-response-map} on the two-dimensional reduced phase space
and does not add an additional physical scalar mode.
